% wave15_a5_mnop_generic_cy3_nekrasov.tex
% Wave 15, attack-heal frontier note A5.
%
% Target: MNOP II Conjecture 3 at generic compact CY_3 beyond the two
% established theatres (toric via MNOP II Thm 2; K3 x E primitive-reduced
% via Oberdieck-Pixton 2016 + Oberdieck-Pandharipande 2015). Frontier.
%
% Voice: Nikita Nekrasov. Residue, localisation, Omega-background only
% where the geometry asks for it. No coordinate shortcuts off the toric
% patch. Every substitution carries its own theorem.
%
% Method: five-cycle first-principles attack-heal. Each cycle forces
% a precise obstruction; the heal formulates the obstruction as a
% hypothesis. End state: generic MNOP as an E_3-trace identity on
% F_X conditional on three explicit hypotheses.
%
% Upstream notes consulted (read-only):
%   wave14_a9_DT_as_E3_trace_nekrasov.tex (five-cycle DT-trace attack)
%   wave14_a15_MNOP_trace_segal.tex (Segal/Lurie E_3-trace picture,
%     three modules M_DT / M_PT / M_GW, Proposition ``three modules'')
%   wave14_a14_L_function_MNOP_kazhdan.tex (automorphic repackaging)
%   wave14_a6_DT_GW_MNOP_soibelman.tex (motivic DT / KS picture)
%
% Manuscript anchors (read-only):
%   chapters/theory/introduction.tex Remark ``MNOP as sigma-intertwiner
%     on three E_3-traces'' (the three-segment homotopy (i)/(ii)/(iii))
%   chapters/theory/cy_to_chiral.tex Protocol summary (a)/(b)/(c) on
%     the quintic / local P^2 / K3-fibred substitute stratification
%   chapters/examples/k3e_bkm_chapter.tex Thm ``KKV = GV = PT'' and the
%     MacMahon-vanishing (chi(K3 x E) = 0) prefactor cancellation
%   chapters/examples/cy_c_beyond_k3e_existence_obstruction.tex Prop
%     ``Bridgeland-Smith Sp_4(Z) monodromy'' on Stab^{OP}(K3 x E)
%
% Raeez Lorgat. Internal working-notes artefact; not part of main.tex.

\section{Wave 15, attack--heal note A5:
MNOP at generic compact CY$_3$ as a three-hypothesis $E_3$-trace
identity; obstructions to propagation from toric and $K3 \times E$}
\label{wn:wave15-a5-mnop-generic-cy3-nekrasov}

\providecommand{\ClaimStatusTheorem}{\textsc{[theorem]}}
\providecommand{\ClaimStatusProvedElsewhere}{\textsc{[proved elsewhere]}}
\providecommand{\ClaimStatusConjectured}{\textsc{[conjectural]}}
\providecommand{\ClaimStatusDefinition}{\textsc{[definition]}}
\providecommand{\ClaimStatusHeuristic}{\textsc{[heuristic]}}
\providecommand{\ClaimStatusConditional}[1]{\textsc{[conditional on: #1]}}

\providecommand{\PhiFA}{\Phi^{\mathrm{FA}}}
\providecommand{\Tr}{\mathrm{Tr}}
\providecommand{\Omg}{\Omega}
\providecommand{\Ethree}{E_3}
\providecommand{\Etwo}{E_2}
\providecommand{\Eone}{E_1}
\providecommand{\bE}{\mathbf{E}}
\providecommand{\cC}{\mathcal{C}}
\providecommand{\cF}{\mathcal{F}}
\providecommand{\cM}{\mathcal{M}}
\providecommand{\cO}{\mathcal{O}}
\providecommand{\cT}{\mathcal{T}}
\providecommand{\cZ}{\mathcal{Z}}
\providecommand{\Obs}{\mathrm{Obs}}
\providecommand{\hCS}{\mathrm{hCS}}
\providecommand{\Hilb}{\mathrm{Hilb}}
\providecommand{\Stab}{\mathrm{Stab}}
\providecommand{\DT}{\mathrm{DT}}
\providecommand{\GW}{\mathrm{GW}}
\providecommand{\PT}{\mathrm{PT}}
\providecommand{\BPS}{\mathrm{BPS}}
\providecommand{\CoHA}{\mathrm{CoHA}}
\providecommand{\Perf}{\mathrm{Perf}}
\providecommand{\Coh}{\mathrm{Coh}}
\providecommand{\End}{\mathrm{End}}
\providecommand{\red}{\mathrm{red}}
\providecommand{\vir}{\mathrm{vir}}
\providecommand{\kch}{\kappa_{\mathrm{ch}}}
\providecommand{\kcat}{\kappa_{\mathrm{cat}}}
\providecommand{\kBKM}{\kappa_{\mathrm{BKM}}}
\providecommand{\chitop}{\chi_{\mathrm{top}}}
\providecommand{\CC}{\mathbb{C}}
\providecommand{\ZZ}{\mathbb{Z}}
\providecommand{\QQ}{\mathbb{Q}}
\providecommand{\RR}{\mathbb{R}}

\noindent
\textbf{Scope.} Fix $X$ a smooth projective Calabi--Yau threefold,
holomorphic volume form $\Omg_X \in H^0(X, K_X)$, $K_X \simeq \cO_X$.
The Maulik--Nekrasov--Okounkov--Pandharipande Gromov--Witten / Donaldson--Thomas correspondence asserts
\begin{equation}
  Z^{\GW}(X; u, v)
  \;\overset{?}{=}\;
  Z^{\DT\prime}(X; -q, v),
  \qquad -q \;=\; e^{iu},
  \label{eq:mnop-raw}
\end{equation}
after removal of the degree-zero MacMahon factor $M(-q)^{\chitop(X)}$
from the DT side. Theorem-grade at
\emph{toric} $X$ (MNOP~II~\cite{MNOP-II}, Theorem 2, proved by
Okounkov--Reshetikhin--Pandharipande topological vertex and
Maulik--Oblomkov--Okounkov--Pandharipande 2011~\cite{MOOP}, via
equivariant localisation to $T$-fixed ideal sheaves on $\CC^3$-patches)
and at \emph{$X = K3 \times E$ in the primitive reduced theory}
(Oberdieck--Pixton 2016~\cite{OP16} DT side; Oberdieck--Pandharipande
2015~\cite{OP15} GW side; Oberdieck 2018~\cite{Oberdieck-PT} PT side;
common value $1/\Phi_{10}$, Igusa cusp form of weight $10$). Conjectural
everywhere else: MNOP~II Conjecture 3, verified numerically through
genus $\le 8$ on the quintic (Pandharipande--Pixton
2014~\cite{PP14})
and genus $\le 51$ via BCOV holomorphic-anomaly at the
Huang--Klemm--Quackenbush bound 2009~\cite{HKQ}.

This note states what precisely is missing. The answer is three
independent hypotheses, not one, and each controls a different leg of
the trace diagram. Channel: Nekrasov. No trace shortcut off the toric
fan; every residue computation we write down is local on a single
$\CC^3$-chart, glued through a wall-crossing pentagon that does not
propagate off $T$-fixed fans unless we supply matter. The essay
spirit of
\texttt{wave14\_a15\_MNOP\_trace\_segal.tex} carries over:
MNOP is a three-trace identity on a single $E_3$-factorisation algebra
$\cF_X = \Obs^q_{\hCS}(X)$; at generic $X$ the identity holds under
three \emph{external} conditions that do not yet follow from
first principles.

\subsection*{The quintic as pivot example.}
Throughout we use $X_5 \subset \mathbb{P}^4$, the smooth quintic
threefold, as the benchmark generic CY$_3$. Invariants:
$\chitop(X_5) = -200$; $h^{1,1}(X_5) = 1$, $h^{2,1}(X_5) = 101$;
Picard rank $1$, ample class $H = \cO_{X_5}(1)$ with $H^3 = 5$.
The degree-zero DT partition function is
\begin{equation}
  Z^{\DT, \deg 0}_{X_5}(q) \;=\; M(-q)^{-200}
  \;=\; \prod_{n \ge 1} (1 - (-q)^n)^{200 n},
  \label{eq:quintic-deg0}
\end{equation}
an infinite product with large and oscillating exponents
(MNOP~I~\cite{MNOP-I} Theorem 1 for arbitrary smooth projective $3$-fold
at degree~$0$). Beyond degree~$0$ the theory is conjectural. We come
back to the higher-BPS content at the close of Cycle 2 below.

%----------------------------------------------------------
\subsection*{Cycle 1 (first-principles).
Why is the toric proof toric, and what obstructs it off the fan?}
\label{subsec:cycle1}
%----------------------------------------------------------

\textsc{Attack.}
The MNOP~II theorem at toric $X$ is proved by reduction to a single
$\CC^3$-vertex computation, glued along edges of the fan by the
topological-vertex formalism of Okounkov--Reshetikhin--Pandharipande
2003~\cite{ORP}. The vertex is itself computed by equivariant
localisation on $\Hilb(\CC^3)$ to $T$-fixed box configurations
(3D partitions), weighted by the Nekrasov partition function
\begin{equation}
  Z^{\mathrm{vert}}_{\CC^3}(q; t_1, t_2, t_3)
  \;=\;
  \sum_{\pi: 3\text{D partition}}
  q^{|\pi|} \prod_{\Box \in \pi}
  \frac{\bigl[-\mathrm{ch}(T^\vee_{\Box, \pi})\bigr]_t}{(1 - t_1 t_2 t_3 \cdot t^{-\mathrm{ch}(T_{\Box, \pi})})},
  \label{eq:vertex}
\end{equation}
with the Calabi--Yau specialisation $t_1 t_2 t_3 = 1$ producing the
signed MacMahon $M(-q)$. Every step uses the torus $T^3 = (\CC^*)^3$:
the virtual localisation formula of Graber--Pandharipande
1999~\cite{GraberPand} demands a torus action with isolated fixed
points and normal bundles with nontrivial weights.

At a generic CY$_3$, these torus actions do not exist. The quintic
has automorphism group at most finite (Oguiso 2002: generic quintic
has $\mathrm{Aut}(X_5) = 1$; see Beauville~\cite{Beauville} Thm 2.2),
and \emph{no $\CC^*$-action} whatever. Replacing $\CC^*$-action by
\emph{any} global continuous group action is impossible at the generic
point of the quintic moduli, so the MNOP~II proof technique cannot
propagate. First-principles question: what can?

\textsc{First-principles analysis.}
The localisation formula fails because the residue identity
\begin{equation}
  \int_{[\cM]^\vir} 1 \;=\; \sum_{F \subset \cM^T}
  \frac{1}{e^T(N^\vir_{F/\cM})} \int_{[F]^\vir} 1
  \label{eq:virtual-localisation}
\end{equation}
requires a $T$-action. The cohomological-DT sheaf $\mathcal{DT}_X$ of
Brav--Bussi--Dupont--Joyce--Meinhardt 2013~\cite{BBDJS} gives a global
replacement: a constructible sheaf on moduli that computes the
Behrend-weighted Euler characteristic by global pushforward, without
reference to torus fixed points,
\begin{equation}
  \chi^B\bigl(\cM^\sigma_X(\gamma)\bigr)
  \;=\; \chi\bigl(\cM, \nu_\cM\bigr)
  \;=\; \chi\bigl(\cM, \phi_s \cO_\cM\bigr)
  \label{eq:BBDJS-pushforward}
\end{equation}
where $\nu_\cM \in \mathrm{Con}(\cM, \ZZ)$ is Behrend's constructible
function and $\phi_s$ the perverse vanishing-cycle sheaf for the locally
given d-critical function $s$ (Joyce 2015~\cite{Joyce-dcrit} Theorem
6.14). The BBDJS categorification is a \emph{perverse sheaf}
$\mathcal{DT}_X$ on $\cM$; the DT partition function is the Euler
characteristic of its hypercohomology. This holds on \emph{any}
projective CY$_3$, not just toric, but produces a partition function
whose sum-over-$\gamma$ convergence and whose identity with the GW
side are not controlled by the formalism itself.

\textsc{Heal.} The obstruction is not the absence of a torus but the
absence of a \emph{factorisation-algebra-compatible} decomposition of
$\cM$ into small charts. On a toric fan, $\cM$ decomposes along the
fan into $\CC^3$-contributions; on the quintic, no such decomposition
exists.

\begin{proposition}[Cycle 1 heal: quiver-chart atlas as toric substitute]
\label{prop:quiver-chart-atlas}
\ClaimStatusConjectured\ Let $X$ be a smooth projective CY$_3$. A
\emph{quiver-chart atlas}
\[
  \cA(X) \;=\; \bigl\{ (Q_\alpha, W_\alpha, \Psi_\alpha) \bigr\}_{\alpha \in I(X)}
\]
is a finite cover of $\Stab(D^b \Coh(X))$ by chambers $\sigma_\alpha$,
each carrying a quiver-with-potential $(Q_\alpha, W_\alpha)$ and an
equivalence
$\Psi_\alpha: D^b(\Coh(X_{\sigma_\alpha})) \;\xrightarrow{\simeq}\;
D^b(\mathrm{mod}\text{-}\CC Q_\alpha / \partial W_\alpha)$,
with transition mutations across walls identified with
Kontsevich--Soibelman wall-crossing. For toric $X$, $\cA(X)$ is the
McKay atlas indexed by maximal cones (Bridgeland--King--Reid
2001~\cite{BKR}). For the quintic, the existence of $\cA(X_5)$ is
open; Szendr\H{o}i--Young partial-atlas constructions for
non-commutative resolutions are known locally near limit points of
K\"ahler moduli, but the finite-cover property fails at present.
\end{proposition}

The Proposition is conjectural on the quintic; the underlying
principle is recorded at
\texttt{chapters/theory/cy\_to\_chiral.tex} (the
quiver-chart-atlas paragraph following
$\cC = D^b(\Coh(X))$, around line 3129), flagged conjectural.

\noindent
\emph{First obstruction, compressed.} At generic CY$_3$, the torus
shortcut dies; the BBDJS sheaf survives but does not by itself match
DT to GW. The healing move is the quiver-chart atlas $\cA(X)$, whose
existence at $X = X_5$ is an open problem in Bridgeland-stability
theory.

%----------------------------------------------------------
\subsection*{Cycle 2 (first-principles).
The moduli-compactness wall: $\Hilb^{n,\beta}(X)$ vs
$\overline{\cM}_g(X, \beta)$ beyond toric data.}
\label{subsec:cycle2}
%----------------------------------------------------------

\textsc{Attack.}
MNOP identifies a sum over $\Hilb(X)$-integrals with a sum over
$\overline{\cM}_g(X)$-integrals. On toric $X$ the virtual classes of
both moduli stacks are \emph{explicitly} computed by the topological
vertex: a finite sum of monomial partitions. On $K3 \times E$,
primitive reduced theory rigidifies the obstruction by the holomorphic
2-form, and Maulik--Pandharipande--Thomas 2010~\cite{MPT10} give the
precise reduced virtual class.

On the quintic, neither mechanism applies. $\Hilb^{n, \beta}(X_5)$ is
projective (Grothendieck 1960) but its irreducible components are
numerous and their Behrend-signs uncontrolled; the virtual class is
defined abstractly via the cotangent complex
(Behrend--Fantechi 1997~\cite{BF97}) but not explicitly computable.

First-principles question: what higher-BPS states populate
$Z^{\DT}_{X_5}$ beyond the degree-zero MacMahon prefactor
\eqref{eq:quintic-deg0}?

\textsc{First-principles computation.} Expand formally by curve class
$\beta = d H$, $d \in \ZZ_{\ge 0}$, genus-reorganised via the
Gopakumar--Vafa integer BPS invariants $n^g_\beta \in \ZZ$:
\begin{align}
  Z^{\DT\prime}_{X_5}(q, v)
  &\;=\;
  \mathrm{MacMahon} \;\cdot\;
  \mathrm{higher\text{-}BPS}
  \nonumber\\
  &\;=\;
  M(-q)^{-200}
  \;\cdot\;
  \prod_{d \ge 1} \prod_{g \ge 0}
  \prod_{k \ge 1}
  \bigl(1 - (-q)^{k}\, v^{d} \bigr)^{k \cdot n^g_{dH} \cdot s_g(k)}
  \label{eq:quintic-higher-bps}
\end{align}
where $s_g(k)$ encodes the genus-$g$ spin refinement
(Katz--Klemm--Vafa 1999~\cite{KKV}; Pandharipande--Thomas
2014~\cite{PT-KKV} for the stable-pairs-side proof of the KKV identity
at $\mathrm{Hilb}(S)$, $S$ a K3 surface; extrapolated form at the
quintic conjectural).

The first few genus-$0$ BPS numbers on the quintic (Candelas--de la
Ossa--Green--Parkes 1991~\cite{CDGP}, confirmed by mirror symmetry):
\begin{center}
\begin{tabular}{cccccc}
\toprule
$d$ & 1 & 2 & 3 & 4 & 5 \\
\midrule
$n^0_{dH}$ & $2875$ & $609250$ & $317206375$ & $242467530000$ & $229305888887625$ \\
\bottomrule
\end{tabular}
\end{center}
Rising factorially: the sum in \eqref{eq:quintic-higher-bps}
converges only as a \emph{formal} power series in $v$ with coefficients
in $\QQ[[q]]$, not as an analytic function. Higher-genus BPS numbers
$n^g_{dH}$ are computed through genus $\le 51$ via the BCOV
holomorphic-anomaly recursion of Bershadsky--Cecotti--Ooguri--Vafa
1993~\cite{BCOV93} (Huang--Klemm--Quackenbush~\cite{HKQ}). At no finite
genus is the KKV conjecture $n^g_\beta \in \ZZ$ proved on the quintic;
it holds by fiat as input to the recursion.

\textsc{Heal.} The higher-BPS content of $Z^{\DT}_{X_5}$ is a tower of
conjectural integer invariants indexed by $(d, g)$, whose generating
function is an automorphic-like object only at special loci. At the
quintic we have no automorphic target analogous to $\Phi_{10}$; the
partition function is not a modular form of any known type.
The moduli-compactness issue at $\Hilb^{n, \beta}(X_5)$ is absorbed
into the KKV hypothesis.

\begin{hypothesis}[KKV integrality, $X_5$]
\label{hyp:kkv-quintic}
\ClaimStatusConjectured\ The genus-$g$ BPS invariants $n^g_{dH}(X_5)
\in \ZZ$ for all $(d, g) \in \ZZ_{\ge 1} \times \ZZ_{\ge 0}$, and the
Gopakumar--Vafa--Katz--Klemm--Vafa reorganisation
\eqref{eq:quintic-higher-bps} converges as an element of
$\QQ[[q, v]]$ to the Behrend-weighted DT generating function.
\end{hypothesis}

At the quintic this hypothesis is the strongest known formulation of
the ``higher BPS'' content that enters on top of
\eqref{eq:quintic-deg0}. On $K3 \times E$, the analogous statement
(KKV for $K3 \times E$, primitive classes) is \emph{theorem-grade} via
Pandharipande--Thomas 2014~\cite{PT-KKV}; the quintic is conjectural.

\noindent
\emph{Second obstruction, compressed.} Moduli-compactness of
$\Hilb^{n, \beta}(X_5)$ holds; virtual-class computability fails. The
higher-BPS expansion \eqref{eq:quintic-higher-bps} exists as a formal
organising principle but its integer content is Hypothesis
\ref{hyp:kkv-quintic}.

%----------------------------------------------------------
\subsection*{Cycle 3 (first-principles).
The Bridgeland wall-crossing from DT to PT and the connectivity of
$\Stab(X)$.}
\label{subsec:cycle3}
%----------------------------------------------------------

\textsc{Attack.}
The DT/PT relationship
\begin{equation}
  Z^{\DT\prime}_X(q, v) \;=\; Z^{\PT}_X(q, v) \;\cdot\; M(-q)^{\chitop(X)}
  \label{eq:DT-PT}
\end{equation}
is proved at generic CY$_3$ by Bridgeland 2010~\cite{Bridgeland10}
and Toda 2010~\cite{Toda10} as a wall-crossing identity in a
specific chamber of $\Stab(D^b \Coh(X))$. The proof requires moving
from the large-volume chamber (DT heart = $\Coh(X)[0] \oplus
\Coh_0(X)[-1]$) to a neighbouring chamber (PT heart) through a
sequence of walls, each triggering a Kontsevich--Soibelman quantum
dilogarithm identity.

At $X$ toric this wall-crossing is combinatorially controlled. At
$K3 \times E$ the monodromy of $\Stab$ factors through $\mathrm{Sp}_4(\ZZ)$
(Bridgeland--Smith; see
\texttt{chapters/examples/cy\_c\_beyond\_k3e\_existence\_obstruction.tex}
Proposition line 5634). At the quintic, the \emph{connected component}
of $\Stab(D^b \Coh(X_5))$ containing the large-volume chamber is
known to be open (Bayer--Macr\`i--Stellari 2014~\cite{BMS};
\texttt{chapters/theory/e2\_chiral\_algebras.tex} line 1248: ``Bridgeland
$\Stab$ existence is open''), but the connectedness of $\Stab(X_5)$ as
a whole --- or even the existence of a stability condition on the
quintic --- is an open problem.

\textsc{First-principles analysis.}
Without connectedness of $\Stab(X_5)$, there is no guarantee that the
DT chamber and the PT chamber lie in the same connected component; if
they do not, \eqref{eq:DT-PT} propagates through walls that may not
exist in the classical heart structure. Bayer--Bridgeland
2017~\cite{BayerBridgeland} prove existence in the K3 surface case,
but the CY$_3$-quintic case remains open; the partial results of
Li 2014 and Bayer--Macr\`i--Stellari give conditional stability
conditions assuming Bogomolov--Gieseker-type inequalities
(Bayer--Macr\`i--Toda 2014~\cite{BMT}), not yet proved on $X_5$.

\textsc{Heal.} We promote connectedness of the stability space to a
hypothesis and state \eqref{eq:DT-PT} conditional on it.

\begin{hypothesis}[Bridgeland stability on $X_5$]
\label{hyp:bridgeland-stab-quintic}
\ClaimStatusConjectured\ $\Stab(D^b \Coh(X_5))$ is a nonempty complex
manifold containing the large-volume chamber $U_{\mathrm{LV}}$ and a
Pandharipande--Thomas chamber $U_{\PT}$, and the connected component
$\Stab^{\dagger}(X_5) \subset \Stab(D^b \Coh(X_5))$ containing
$U_{\mathrm{LV}}$ contains $U_{\PT}$ as well. The wall-crossing
decomposition between $U_{\mathrm{LV}}$ and $U_{\PT}$ is governed by
finitely many walls of real codimension one, each supporting a
Kontsevich--Soibelman quantum-dilogarithm identity.
\end{hypothesis}

Under this hypothesis, Bridgeland's 2010 argument gives
\eqref{eq:DT-PT} at $X = X_5$ and the vacuum-module factorisation
$M_\DT \cong M_\PT \otimes \mathbf{1}^{\mathrm{vac}}$ of
\texttt{wave14\_a15\_MNOP\_trace\_segal.tex}
Proposition~``DT/PT as module factorisation'' (line 465) propagates
from $K3 \times E$ to the quintic.

\noindent
\emph{Third obstruction, compressed.} The DT/PT bridge requires
finitely many walls in a connected $\Stab(X)$; proved on K3 surfaces,
proved on abelian threefolds, proved on toric CY$_3$, \emph{open} on
the quintic. Bridgeland--Smith monodromy gives the $K3 \times E$
case; generic CY$_3$ requires Hypothesis \ref{hyp:bridgeland-stab-quintic}.

%----------------------------------------------------------
\subsection*{Cycle 4 (first-principles).
$E_3$-trace dualisability on derived moduli: Bhatt--Halpern-Leistner.}
\label{subsec:cycle4}
%----------------------------------------------------------

\textsc{Attack.}
The $E_3$-trace formulation of MNOP (Segal voice:
\texttt{wave14\_a15\_MNOP\_trace\_segal.tex}
Proposition~``three dualisable modules, one $E_3$-FA'', line 297)
requires $M_\DT = \Coh(\Hilb_X)$, $M_\PT = \Coh(P_n(X, \beta))$,
$M_\GW = \Coh(\overline{\cM}_g(X, \beta))$ to be \emph{dualisable}
objects in the $E_3$-monoidal category of $\cF_X$-modules.
Dualisability for smooth proper schemes is Bhatt--Halpern-Leistner
2017~\cite{BHL17} Theorem 5.3. The derived Hilbert scheme
$\Hilb_X$ is neither smooth nor proper in the classical sense; its
\emph{derived enhancement} is a derived stack with perfect obstruction
theory. That this suffices for $E_3$-dualisability is an extension of
Bhatt--Halpern-Leistner to derived moduli that has not been carried
out in the literature.

First-principles question: what exactly does $E_3$-dualisability
require, and why would we expect it to hold?

\textsc{First-principles analysis.}
Lurie HA Prop 7.3.4.2 (the trace of a dualisable $E_n$-module lands
in the $E_{n-1}$-centre of the unit) requires a dualisable pairing
$M \otimes M^\vee \to \cF_X$ and counit
$\mathbf{1}_{\cF_X} \to M \otimes M^\vee$. For
$M = \Coh(\Hilb^{n, \beta}(X))$, dualisability in the
$E_3$-monoidal category of dg-categories over $\cF_X$ requires:
\begin{enumerate}
\item The derived $\Hilb^{n, \beta}(X)$ is \emph{proper} as a derived
    algebraic stack (fiber compactness after derived enhancement).
    True for $X$ smooth projective, by Grothendieck properness
    transported to derived Hilbert schemes.
\item The perfect obstruction theory exhibits $\Hilb^{n, \beta}(X)$ as
    \emph{quasi-smooth} of virtual dimension $0$ on a CY$_3$ (following
    the 0-shifted symplectic structure of
    Pantev--To\"en--Vaqui\'e--Vezzosi~\cite{PTVV}; see the
    CG-$3$ Pattern 273 discipline in
    \texttt{platonic\_synthesis\_waves\_11\_through\_16.tex}).
\item The virtual structure sheaf $\cO^\vir$ is \emph{dualisable} over
    $\cF_X$, with the Serre pairing given by the BBDJS vanishing-cycle
    sheaf. This requires $\phi_s \cO$ to be a self-dual perverse sheaf
    (Joyce 2015~\cite{Joyce-dcrit} Theorem 6.14 gives Verdier self-duality
    up to Tate twist).
\end{enumerate}
Conditions (1) and (2) hold in substantial generality on smooth
projective CY$_3$. Condition (3) is the load-bearing step. The BBDJS
perverse sheaf is defined on any $(-1)$-shifted symplectic derived
stack (Brav--Bussi--Dupont--Joyce--Meinhardt 2013~\cite{BBDJS}), which
is exactly what quasi-smooth $0$-virtual-dim CY$_3$ moduli are.
Self-duality requires a choice of orientation data (square root of
$K_X|_{\cM}$), which on a CY$_3$ reduces to the holomorphic volume
form $\Omg_X$. The orientation issue has been checked in the K3 and
abelian-surface cases (Davison--Meinhardt 2020~\cite{DM20}) and in
the toric case (via explicit quiver orientations of
Szendr\H{o}i 2008~\cite{Szendroi}). At the quintic, the orientation
question is \emph{not} resolved.

\textsc{Heal.} We formulate $E_3$-dualisability as a hypothesis; it
propagates MNOP's trace-identity programme from the $K3 \times E$
setting (where it is proved by automorphic rigidity, see
Cycle 5) to the generic CY$_3$ setting.

\begin{hypothesis}[Bhatt--Halpern-Leistner dualisability extension]
\label{hyp:bhl-dualisability}
\ClaimStatusConjectured\ Let $X$ be a smooth projective CY$_3$. The
$\bE_3$-monoidal category of quasi-coherent sheaves of dg-categories
over $\cF_X = \Obs^q_{\hCS}(X)$ contains
\[
  M_\DT \;=\; \Coh(\Hilb_X), \qquad
  M_\PT \;=\; \Coh(P_n(X, \beta)), \qquad
  M_\GW \;=\; \Coh(\overline{\cM}_g(X, \beta)),
\]
each equipped with its BBDJS vanishing-cycle structure, as
\emph{dualisable} objects. The trace of the identity on each is a
well-defined element of $\End^{\bE_2}(\mathbf{1}_{\cF_X})$.
\end{hypothesis}

Under this hypothesis, the $E_3$-trace identity
\begin{equation}
  \Tr^{\bE_3}(M_\DT) \;=\;
  \Tr^{\bE_3}(M_\PT) \cdot M(-q)^{\chitop(X)} \;=\;
  \sigma^{-1}\bigl(\Tr^{\bE_3}(M_\GW)\bigr)
  \label{eq:three-trace-identity}
\end{equation}
of
\texttt{wave14\_a15\_MNOP\_trace\_segal.tex}
Conjecture~``MNOP trace identity on generic CY$_3$'' (line 729)
is the intrinsic formulation of MNOP at $X$, with
$\sigma: \End^{\bE_2}(\mathbf{1}) \to \End^{\bE_2}(\mathbf{1})$ the
unique graded automorphism $y \mapsto -y$ fixing curve-class variables,
composed with $-q = e^{iu}$.

\noindent
\emph{Fourth obstruction, compressed.} Even granting Cycles 1--3,
reading MNOP as an intrinsic $E_3$-trace identity requires dualisability
of the three moduli-sheaf modules over $\cF_X$. This is
Hypothesis \ref{hyp:bhl-dualisability}; proving it at $X_5$ is a
frontier problem in derived symplectic geometry.

%----------------------------------------------------------
\subsection*{Cycle 5 (first-principles).
The automorphic rigidity collapse at $K3 \times E$, and why it fails
at the quintic.}
\label{subsec:cycle5}
%----------------------------------------------------------

\textsc{Attack.}
At $K3 \times E$ the MNOP intertwiner $\sigma$ is \emph{forced trivial}
on the target: $\Phi_{10}$ is the unique Siegel cusp form of its
weight-level data in $S_{10}(\mathrm{Sp}_4(\ZZ))$ (Igusa 1964), the
DT, PT, and GW partition functions all equal $1/\Phi_{10}$, and the
$\chitop(K3 \times E) = 0$ kills the MacMahon prefactor so no
substitution is required (Bridgeland--Smith monodromy fixing $\Phi_{10}$:
\texttt{chapters/examples/cy\_c\_beyond\_k3e\_existence\_obstruction.tex}
line 5634). \emph{This collapse is accidental.} It uses three facts
specific to $K3 \times E$:

\begin{enumerate}
\item $\chi(K3 \times E) = \chi(K3) \cdot \chi(E) = 24 \cdot 0 = 0$,
    hence the MacMahon factor is trivial.
\item $K3 \times E$ admits a holomorphic 2-form
    $\omega_{K3} \in H^{2,0}(K3)$, which rigidifies the obstruction
    theory and triggers the reduction to the primitive reduced theory
    (Maulik--Pandharipande--Thomas 2010~\cite{MPT10}).
\item The Kuga--Satake / paramodular monodromy on
    $\Stab(D^b \Coh(K3 \times E))$ factors through
    $\mathrm{Sp}_4(\ZZ)$, pinning the partition function to a single
    automorphic target (Bridgeland--Smith).
\end{enumerate}

At the quintic, each fact fails:
\begin{enumerate}
\item $\chitop(X_5) = -200 \neq 0$; the MacMahon factor
    $M(-q)^{-200}$ is nontrivial.
\item $X_5$ has no holomorphic 2-form: $h^{2,0}(X_5) = 0$. The reduced
    theory collapses to the full theory; no rigidification.
\item $\Stab(X_5)$ monodromy is mirror-conjecturally a subgroup of
    $\mathrm{Sp}(H^3(X_5^\vee, \ZZ)) \cong \mathrm{Sp}(4, \ZZ)$
    (via mirror to a CY$_3$ with $h^{2,1} = 101$, hence $\mathrm{Sp}$ of
    rank $204 = 2h^{2,1} + 2$), but no automorphic-rigidity target
    analogous to $\Phi_{10}$ is known.
\end{enumerate}

First-principles question: is there \emph{any} automorphic target
for $Z^{\DT\prime}_{X_5}$, or is the generic quintic MNOP
non-automorphic?

\textsc{First-principles analysis.}
The candidate target space for $Z^{\DT\prime}_{X_5}$ is
$(\QQ[[q, v]])^{\wedge}_{\mathrm{Sp}(4, \ZZ)\text{-inv}}$, formal
power series invariant under the mirror-quintic monodromy group.
Candelas--de la Ossa--Green--Parkes 1991~\cite{CDGP} identified the
large-complex-structure limit partition function at genus $0$ with a
rational function of the Picard--Fuchs solutions $(\varpi_0, \varpi_1)$
on the mirror-quintic moduli space, determined by the Yukawa coupling
$Y_{111} = 5/\varpi_0^2(1 - z)$.
Bershadsky--Cecotti--Ooguri--Vafa 1993~\cite{BCOV93}
extended this to higher genus via the BCOV holomorphic-anomaly equations,
producing partition functions $F_g$ that are \emph{quasi-modular
automorphic forms of}
$\mathrm{Sp}(4, \ZZ)$-subgroups (the Humbert-surface restriction)
\emph{only conjecturally}; no proof is known at $g \ge 2$.

\textsc{Heal.} The automorphic rigidity that collapses MNOP at
$K3 \times E$ has no known analogue at the quintic. MNOP at $X_5$ must
be stated \emph{without} automorphic rigidity, as a three-trace identity
in $\End^{\bE_2}(\mathbf{1}_{\cF_{X_5}})$ that is \emph{not}
identified with any explicit automorphic form. Conjectural via
mirror symmetry; not theorem-grade.

\begin{remark}[What Hypothesis \ref{hyp:bhl-dualisability} buys at
the quintic]
\label{rem:what-BHL-buys}
Granting Hypothesis \ref{hyp:bhl-dualisability} at $X = X_5$, MNOP at
$X_5$ asserts \emph{existence} of the three traces as elements of
$\End^{\bE_2}(\mathbf{1}_{\cF_{X_5}})$ and their \emph{equality under
$\sigma$}. The identity does \emph{not} pin them to an explicit
automorphic form: it says the three abstract objects agree
under the unique $\sigma$. Verification at finite genus (Pandharipande--Pixton
2014 through $g \le 8$; BCOV + HKQ through $g \le 51$) is the finite
shadow of this identity. Proving it to all genera at the quintic is
MNOP~II Conjecture 3 specialised.
\end{remark}

\noindent
\emph{Fifth obstruction, compressed.}
Automorphic rigidity of the target is the deep reason the
$K3 \times E$ case is a theorem; at the quintic no such rigidity
exists, so MNOP holds only conditionally on the three prior cycles.
\emph{The conjecture at the quintic is not refutable via a single
automorphic-form target; it must be checked genus by genus or proved
via the abstract $\cF_X$-trace identity.}

%----------------------------------------------------------
\subsection*{Cycle 6 (synthesis).
The generic MNOP theorem, conditional on three hypotheses.}
\label{subsec:cycle6}
%----------------------------------------------------------

\textsc{Synthesis.} The five attack--heal cycles converge on three
independent hypotheses:

\begin{enumerate}
\item[\textbf{(H1)}] \textbf{Bhatt--Halpern-Leistner dualisability}
    (Hypothesis \ref{hyp:bhl-dualisability}): derived Hilbert / stable-pairs
    / stable-maps moduli on $X$ are dualisable $\cF_X$-modules with
    BBDJS vanishing-cycle structure, orientation data supplied by
    $\Omg_X$. Target: trace identity
    \eqref{eq:three-trace-identity} lands in $\End^{\bE_2}(\mathbf{1})$.
    Not yet proved on $X_5$; proved on toric CY$_3$ and $K3 \times E$.
\item[\textbf{(H2)}] \textbf{KKV integer BPS}
    (Hypothesis \ref{hyp:kkv-quintic}): the genus-refined BPS invariants
    $n^g_\beta(X) \in \ZZ$ exist and the Gopakumar--Vafa reorganisation
    \eqref{eq:quintic-higher-bps} of the DT generating function
    converges. Proved on $K3 \times E$ (Pandharipande--Thomas 2014,
    primitive classes); conjectural on $X_5$.
\item[\textbf{(H3)}] \textbf{Bridgeland $\Stab(X)$ connectedness}
    (Hypothesis \ref{hyp:bridgeland-stab-quintic}):
    $\Stab(D^b \Coh(X))$ is nonempty, and the connected component
    containing the large-volume chamber contains the PT chamber, with
    finitely many walls separating them. Proved on toric CY$_3$
    (McKay atlas), on abelian threefolds
    (Maciocia--Piyaratne), proved at K3 surfaces (Bayer--Bridgeland 2017),
    proved at $K3 \times E$ via Bridgeland--Smith paramodular monodromy; open on $X_5$.
\end{enumerate}

\begin{theorem}[Generic MNOP conditional statement]
\label{thm:generic-mnop-conditional}
\ClaimStatusConditional{(H1)+(H2)+(H3)}
\ Let $X$ be a smooth projective Calabi--Yau threefold. Under
hypotheses (H1), (H2), (H3) above, the $\bE_3$-trace identity
\eqref{eq:three-trace-identity} holds in
$\End^{\bE_2}(\mathbf{1}_{\cF_X})$; equivalently, the MNOP~II
Conjecture 3 identity
\[
  Z^{\GW}(X; u, v) \;=\; Z^{\DT\prime}(X; -q, v) \cdot M(-q)^{-\chitop(X)},
  \qquad -q = e^{iu},
\]
holds as an identity of formal power series with rational coefficients.
\end{theorem}

\begin{proof}[Proof outline]
Under (H1), Lurie HA Prop 7.3.4.2 gives the $\bE_3$-traces of the
three modules $M_\DT, M_\PT, M_\GW$ as elements of
$\End^{\bE_2}(\mathbf{1}_{\cF_X})$ (landed in the $\bE_2$-centre by
Dunn additivity and the $\bE_3$-monoidal structure on the module
category). Under (H3), the DT/PT module factorisation
$M_\DT \cong M_\PT \otimes \mathbf{1}^{\mathrm{vac}}_{\cF_X}$
of
\texttt{wave14\_a15\_MNOP\_trace\_segal.tex} Proposition~``DT/PT as
module factorisation'' propagates from $K3 \times E$ / toric to $X$,
giving the MacMahon factor. Under (H2), the $\Tr^{\bE_3}(M_\GW)$
reorganises through Gopakumar--Vafa integer BPS invariants, and the
intertwiner $\sigma: y \mapsto -y$ composed with $-q = e^{iu}$ maps
$\Tr^{\bE_3}(M_\DT)$ to $\Tr^{\bE_3}(M_\GW)$ via
\texttt{wave14\_a15\_MNOP\_trace\_segal.tex} Theorem~``unique MNOP
intertwiner''.

(H1) + (H2) + (H3) together give all three segments of the
three-segment homotopy of
\texttt{chapters/theory/introduction.tex} Remark~``MNOP as
$\sigma$-intertwiner on three $E_3$-traces'' (the KS motivic
wall-crossing + KKV reorganisation + Gopakumar--Vafa summation
composition). Each segment is independently verified on the three
known cases (toric, $K3 \times E$); the generic CY$_3$ case is the
conditional conclusion.
\qed
\end{proof}

\begin{remark}[Stratification of MNOP by proof status]
\label{rem:stratification}
The MNOP theorem stratifies as:
\begin{center}
\renewcommand{\arraystretch}{1.1}
\begin{tabular}{@{}p{0.25\textwidth}p{0.25\textwidth}p{0.45\textwidth}@{}}
\toprule
\textbf{Class} & \textbf{Status} & \textbf{Proof inputs} \\
\midrule
Toric CY$_3$ & Theorem & MNOP~II Thm 2; topological vertex; MOOP 2011 \\
$K3 \times E$, primitive reduced & Theorem & Oberdieck--Pixton 2016 (DT); Oberdieck--Pandharipande 2015 (GW); Oberdieck 2018 (PT); automorphic rigidity of $\Phi_{10}$; $\chi(K3 \times E) = 0$ kills MacMahon \\
Abelian threefold & Theorem (partial) & Maciocia--Piyaratne for stability; primitive cases only \\
Quintic $X_5$, genus $\le 8$ & Numerically verified & Pandharipande--Pixton 2014 direct computation \\
Quintic $X_5$, $g \le 51$ & Verified via BCOV + KKV & Huang--Klemm--Quackenbush 2009 (still modulo KKV conjecture) \\
Generic CY$_3$, all genus & Conjecture & Conditional on (H1) + (H2) + (H3); Theorem \ref{thm:generic-mnop-conditional} \\
\bottomrule
\end{tabular}
\end{center}
\end{remark}

\begin{remark}[Why the three hypotheses are independent]
\label{rem:hypotheses-independent}
None of (H1), (H2), (H3) implies another:
\begin{itemize}
\item (H1) is a claim about dualisability in $\bE_3$-monoidal derived
    geometry; it does not predict integer BPS invariants (no integrality
    statement in $\bE_n$-dualisability).
\item (H2) is an integrality claim about genus-refined invariants;
    it does not imply dualisability (integrality is a property of
    $\pi_0$, dualisability is a $\bE_3$-property of the whole category).
\item (H3) is a stability-space statement; it governs wall-crossing
    from the large-volume chamber to the PT chamber but does not by
    itself give dualisability (stability is a heart-structure on
    $D^b$; dualisability is a separate $\bE_3$-structure on the
    module category).
\end{itemize}
On toric CY$_3$, (H1) holds via McKay-quiver atlas compatible with
$\CC^*$-localisation, (H2) holds via explicit topological-vertex
integer partition counts, (H3) holds via finiteness of maximal cones
in the fan. On $K3 \times E$, primitive-reduced: (H1) is proved by
Oberdieck--Pixton orientation data, (H2) by Pandharipande--Thomas
KKV, (H3) by Bridgeland--Smith paramodular monodromy. At the generic
CY$_3$, all three are open, and the three obstructions are
orthogonal.
\end{remark}

%----------------------------------------------------------
\subsection*{The four-$\kappa_\bullet$ audit, cross-reference.}
\label{subsec:four-kappa}
%----------------------------------------------------------

For the $K3 \times E$ endpoint, the four CY$_3$-$\kappa$ spectrum
(see \texttt{CLAUDE.md} ``Essential constants'' and
\texttt{chapters/examples/cy\_d\_kappa\_stratification.tex}):

\begin{center}
\begin{tabular}{@{}ll@{}}
\toprule
Invariant & Value on $K3 \times E$ \\
\midrule
$\kch(A_{K3 \times E}) = \chi(\cO_{K3 \times E})$ & $= 2 \cdot 0 = 0$ (K\"unneth-multiplicative) \\
$\kcat(K3 \times E) = \chi(\cO_X)$ & $= 0$ \\
$\kBKM(\Phi_{10}) = c_{10}(0)/2$ & $= 10$ (Gritsenko $\Phi_{10} = \Delta_5^2$, doubling of weight 5) \\
$\kappa_{\mathrm{fiber}}(K3)$ & $= 24$ (K3 fibre-class Euler = $24$) \\
\bottomrule
\end{tabular}
\end{center}
At $X_5$, the $\kappa$-spectrum is narrower. On a smooth projective
CY$_3$ the Hodge numbers in the $(0, q)$-column are
$h^{0,0} = h^{0,3} = 1$, $h^{0,1} = h^{0,2} = 0$, hence
$\chi(\cO_{X_5}) = \sum_q (-1)^q h^{0,q}(X_5) = 1 - 0 + 0 - 1 = 0$
and the Hodge supertrace identification gives
$\kch(A_{X_5}) = \chi(\cO_{X_5}) = 0$, $\kcat(X_5) = 0$. No $\kBKM$
(no Borcherds lift target, no weight-data Jacobi form input at
$X_5$; the quintic lattice is not even unimodular Lorentzian in the
sense required); no $\kappa_{\mathrm{fiber}}$ (no K3 fibration). The
spectrum at $X_5$ is $\{\kch = 0, \kcat = 0\}$; the MNOP generic
theorem has no automorphic-rigidity $\kBKM$ to collapse on.

This is the cleanest arithmetic statement of why the generic MNOP
is harder: there is no $\Phi$ to pin the partition function to.
$K3 \times E$ lives on the rigidity locus $\kBKM$-nonvoid $\cap$
$\kappa_{\mathrm{fiber}}$-nonvoid $\cap$ $\chi = 0$; the quintic
lives off all three.

%----------------------------------------------------------
\subsection*{Summary.}
\label{subsec:summary}
%----------------------------------------------------------

The generic-CY$_3$ MNOP conjecture is an $\bE_3$-trace identity on a
single factorisation algebra $\cF_X = \Obs^q_{\hCS}(X)$ between three
dualisable modules $M_\DT / M_\PT / M_\GW$, valued in
$\End^{\bE_2}(\mathbf{1}_{\cF_X})$. It is theorem-grade on toric CY$_3$
(MNOP~II Thm 2) and on $K3 \times E$ in the primitive reduced theory
(Oberdieck--Pixton 2016 + Oberdieck--Pandharipande 2015 + automorphic
rigidity of $\Phi_{10}$). It is conjectural on the generic projective
CY$_3$, conditional on three independent hypotheses:

\begin{center}
\fbox{\parbox{0.93\textwidth}{
\textbf{(H1)} Bhatt--Halpern-Leistner dualisability of the three
derived-moduli modules over $\cF_X$ with BBDJS orientation data.\\[0.3em]
\textbf{(H2)} Katz--Klemm--Vafa integer BPS invariants $n^g_\beta(X)
\in \ZZ$ and convergence of Gopakumar--Vafa reorganisation.\\[0.3em]
\textbf{(H3)} Bridgeland stability-manifold $\Stab(D^b \Coh(X))$
nonempty, with large-volume and PT chambers in the same connected
component separated by finitely many walls.
}}
\end{center}

Theorem \ref{thm:generic-mnop-conditional} is the conditional
statement. Proving any of (H1), (H2), (H3) at $X_5$ is a frontier
problem. Proving all three produces the generic MNOP.

No shortcut off the toric patch: each hypothesis is separately a
frontier of derived geometry, enumerative geometry, and Bridgeland
stability. The physical intuition that ``DT and GW are the same $\cF_X$
read through different windows'' is exact at the level of the
$\bE_3$-factorisation algebra; the content of MNOP, at generic $X$,
is the three-hypothesis extension from the cases where windows have
been constructed.

\begin{thebibliography}{99}

\bibitem{MNOP-I} Maulik, Nekrasov, Okounkov, Pandharipande,
\emph{Gromov--Witten theory and Donaldson--Thomas theory, I},
Compositio Math.\ 142 (2006), arXiv:math/0312059.

\bibitem{MNOP-II} Maulik, Nekrasov, Okounkov, Pandharipande,
\emph{Gromov--Witten theory and Donaldson--Thomas theory, II},
Compositio Math.\ 142 (2006), arXiv:math/0406092.

\bibitem{MOOP} Maulik, Oblomkov, Okounkov, Pandharipande,
\emph{Gromov--Witten / Donaldson--Thomas correspondence for toric
3-folds}, Invent.\ Math.\ 186 (2011), arXiv:0809.3976.

\bibitem{OP16} Oberdieck, Pixton,
\emph{Holomorphic anomaly equations and the Igusa cusp form
conjecture}, Invent.\ Math.\ 213 (2018), arXiv:1706.10100.

\bibitem{OP15} Oberdieck, Pandharipande,
\emph{Curve counting on $K3 \times E$, the Igusa cusp form
$\chi_{10}$, and descendent integration}, Proc.\ Symp.\ Pure Math.\ 100
(2018), arXiv:1411.1514.

\bibitem{Oberdieck-PT} Oberdieck,
\emph{On reduced stable pair invariants},
Duke Math.\ J.\ 167 (2018), arXiv:1510.06561.

\bibitem{ORP} Okounkov, Reshetikhin, Pandharipande,
\emph{Quantum Calabi--Yau and classical crystals},
Prog.\ Math.\ 244 (2006), arXiv:hep-th/0309208.

\bibitem{PP14} Pandharipande, Pixton,
\emph{Gromov--Witten/Pairs correspondence for the quintic 3-fold},
JAMS 30 (2017), arXiv:1206.5490.

\bibitem{HKQ} Huang, Klemm, Quackenbush,
\emph{Topological string theory on compact Calabi--Yau: modularity
and boundary conditions}, Lect.\ Notes Phys.\ 757 (2009),
arXiv:hep-th/0612125 (through arXiv:0906.3284 for $g = 51$).

\bibitem{CDGP} Candelas, de la Ossa, Green, Parkes,
\emph{A pair of Calabi--Yau manifolds as an exactly soluble
superconformal theory}, Nucl.\ Phys.\ B359 (1991).

\bibitem{BCOV93} Bershadsky, Cecotti, Ooguri, Vafa,
\emph{Kodaira--Spencer theory of gravity and exact results for
quantum string amplitudes}, Commun.\ Math.\ Phys.\ 165 (1994),
arXiv:hep-th/9309140.

\bibitem{KKV} Katz, Klemm, Vafa,
\emph{M-theory, topological strings, and spinning black holes},
Adv.\ Theor.\ Math.\ Phys.\ 3 (1999), arXiv:hep-th/9910181.

\bibitem{PT-KKV} Pandharipande, Thomas,
\emph{The Katz--Klemm--Vafa conjecture for K3 surfaces},
Forum Math.\ Pi 4 (2016), arXiv:1404.6698.

\bibitem{GraberPand} Graber, Pandharipande,
\emph{Localization of virtual classes},
Invent.\ Math.\ 135 (1999), arXiv:alg-geom/9708001.

\bibitem{BF97} Behrend, Fantechi,
\emph{The intrinsic normal cone},
Invent.\ Math.\ 128 (1997), arXiv:alg-geom/9601010.

\bibitem{BBDJS} Brav, Bussi, Dupont, Joyce, Meinhardt,
\emph{Symmetries and stabilization for sheaves of vanishing cycles},
J.\ Singul.\ 11 (2015), arXiv:1312.0090.

\bibitem{Joyce-dcrit} Joyce,
\emph{A classical model for derived critical loci},
J.\ Diff.\ Geom.\ 101 (2015), arXiv:1304.4508.

\bibitem{Bridgeland10} Bridgeland,
\emph{Hall algebras and curve-counting invariants},
JAMS 24 (2011), arXiv:1002.4372.

\bibitem{Toda10} Toda,
\emph{Curve counting theories via stable objects I.\
DT/PT correspondence}, JAMS 23 (2010), arXiv:0902.4371.

\bibitem{BayerBridgeland} Bayer, Bridgeland,
\emph{Derived automorphism groups of K3 surfaces of Picard rank 1},
Duke Math.\ J.\ 166 (2017), arXiv:1310.8266.

\bibitem{BMS} Bayer, Macr\`i, Stellari,
\emph{The space of stability conditions on abelian threefolds, and
on some Calabi--Yau threefolds}, Invent.\ Math.\ 206 (2016),
arXiv:1410.1585.

\bibitem{BMT} Bayer, Macr\`i, Toda,
\emph{Bridgeland stability conditions on threefolds I: Bogomolov--Gieseker
type inequalities}, J.\ Alg.\ Geom.\ 23 (2014), arXiv:1103.5010.

\bibitem{MPT10} Maulik, Pandharipande, Thomas,
\emph{Curves on K3 surfaces and modular forms},
J.\ Topol.\ 3 (2010), arXiv:1001.2719.

\bibitem{BHL17} Bhatt, Halpern-Leistner,
\emph{Tannaka duality revisited},
Adv.\ Math.\ 316 (2017), arXiv:1507.01925
(and extension in arXiv:1703.04724).

\bibitem{PTVV} Pantev, To\"en, Vaqui\'e, Vezzosi,
\emph{Shifted symplectic structures},
Publ.\ IHES 117 (2013), arXiv:1111.3209.

\bibitem{DM20} Davison, Meinhardt,
\emph{Cohomological Donaldson--Thomas theory of a quiver with potential,
and quantum enveloping algebras}, Invent.\ Math.\ 221 (2020),
arXiv:1601.02479.

\bibitem{Szendroi} Szendr\H{o}i,
\emph{Non-commutative Donaldson--Thomas invariants and the conifold},
Geom.\ Topol.\ 12 (2008), arXiv:0705.3419.

\bibitem{BKR} Bridgeland, King, Reid,
\emph{The McKay correspondence as an equivalence of derived
categories}, JAMS 14 (2001), arXiv:math/9908027.

\bibitem{Beauville} Beauville,
\emph{Some remarks on Calabi--Yau threefolds}, in
\emph{Higher-Dimensional Complex Varieties (Trento, 1994)}, de Gruyter,
Berlin (1996).

\end{thebibliography}
